\chapter{La qualità del software}

Nei capitoli precedenti la parola \textit{qualità} è comparsa continuamente: l'Ingegneria del Software vuole produrre
«software di alta qualità, nei tempi e nei costi previsti». Ma cosa significa, esattamente, che un software è di qualità?

Sapere rispondere a questa domanda serve da due punti di vista opposti:
\begin{itemize}
  \item come \textbf{produttori}, per capire se ciò che stiamo costruendo è davvero di qualità (e per dimostrarlo al cliente);
  \item come \textbf{committenti o finanziatori}, per capire se il software che stiamo pagando vale quello che costa.
\end{itemize}

Il problema è che la qualità è \textbf{difficile da quantificare}: non esiste un unico numero che la riassuma. Per questo servono
più definizioni, organizzate in modelli condivisi.

% ============================================================
\section{Cos'è la qualità di un prodotto}

La qualità di un prodotto è un concetto \textbf{complesso}. Proviamo a pensarci fuori dal software: cosa rende «di qualità»
questi oggetti?

\begin{center}
\begin{tikzpicture}
  \foreach \i/\titolo/\col/\items in {
    0/{Automobile}/swblue/{Si guasta raramente?\\Consuma poco carburante?\\Costa poco mantenerla?},
    1/{Orologio}/sworange/{È molto preciso?\\Resiste ad acqua e polvere?\\Resiste ai graffi?},
    2/{Componenti meccanici}/swteal/{Hanno tolleranze ridotte?\\Rispettano le specifiche?\\Resistono all'usura?}}{
    \node[chip=\col, minimum width=5.3cm, minimum height=0.75cm, font=\small\bfseries] (h\i) at (\i*5.8,0) {\titolo};
    \node[card=\col, minimum width=5.3cm, text width=4.9cm, align=left, font=\small, below=0.1cm of h\i] {\items};
  }
\end{tikzpicture}
\end{center}

Già da questi esempi emergono due cose:
\begin{itemize}
  \item la qualità dipende dal \textbf{tipo di prodotto}: la precisione conta per un orologio, i consumi per un'auto;
  \item la qualità dipende da \textbf{chi la giudica}: per il guidatore conta il consumo, per chi la ripara la facilità di manutenzione.
\end{itemize}

% ============================================================
\section{Cos'è la qualità del software?}

Con il software è ancora più difficile. Vediamo tre esempi famosi.

\subsection{Esempio 1: efficiente, ma comprensibile?}

Questa è l'implementazione della radice quadrata inversa ($1/\sqrt{x}$) contenuta nel videogioco \textit{Quake III Arena},
attribuita a John Carmack. Il calcolo di $1/\sqrt{x}$ è usatissimo in computer grafica (ad esempio per calcolare angoli di
incidenza e riflessione della luce).

\Needspace{22\baselineskip}\begin{lstlisting}[language=C]
float Q_rsqrt(float number) {
    long i;
    float x2, y;
    const float threehalfs = 1.5F;

    x2 = number * 0.5F;
    y  = number;
    i  = * ( long * ) &y;
    i  = 0x5f3759df - ( i >> 1 );
    y  = * ( float * ) &i;
    y  = y * ( threehalfs - ( x2 * y * y ) );

    return y;
}
\end{lstlisting}

All'epoca questa approssimazione era circa \textbf{4 volte più veloce} di \texttt{(float)(1.0/sqrt(x))} e molti la considerano
un esempio di programmazione geniale. Ma sapresti dire come funziona?

\info{Come funziona (in breve) e perché è un problema}{
  Il codice legge i bit del \texttt{float} come se fossero un intero, li combina con la «costante magica» \texttt{0x5f3759df}
  per ottenere una prima stima del risultato, riconverte i bit in \texttt{float} e infine raffina la stima con un passo del
  metodo di Newton. Il risultato è velocissimo, ma:
  \begin{itemize}
    \item è quasi \textbf{incomprensibile} senza una spiegazione matematica dettagliata (la costante non ha nome né commento);
    \item è \textbf{fragile}: presuppone che \texttt{long} occupi 32 bit come il \texttt{float}; su molti sistemi a 64 bit
          \texttt{long} occupa 64 bit e il codice non funziona più correttamente;
    \item è difficile da \textbf{modificare} e da \textbf{testare} in modo affidabile.
  \end{itemize}
  Quindi: essere efficienti è sinonimo di alta qualità? \textbf{Dipende da quale qualità ci interessa.}
}

\subsection{Esempio 2 e 3: sicuro? facile da usare?}

\begin{center}
\begin{tikzpicture}
  \node[card=swred, text width=8cm, align=left, minimum height=3.6cm] (obsd) {%
    \textbf{\large OpenBSD}\\[3pt]\small
    Molti lo considerano un sistema operativo di altissima qualità: è molto stabile e difficile da violare.
    Secondo il progetto stesso, nell'installazione predefinita sono state scoperte solo \textbf{due} vulnerabilità
    sfruttabili da remoto in tutta la sua storia.\\[4pt]
    \textit{Essere molto sicuri contro gli attacchi significa essere di alta qualità?}};
  \node[card=swgreen, text width=8cm, align=left, minimum height=3.6cm, right=0.5cm of obsd] (vim) {%
    \textbf{\large Vim}\\[3pt]\small
    Molti lo considerano il miglior editor di testo in assoluto: gli utenti esperti lavorano a velocità impressionanti.
    Altri lo trovano troppo difficile da imparare (è celebre la domanda «come si esce da Vim?») e quindi di bassa qualità.\\[4pt]
    \textit{La qualità dipende da chi lo usa?}};
\end{tikzpicture}
\end{center}

\subsection{La lezione: tanti punti di vista}

Non esiste \textbf{una sola} qualità, ma molti approcci e punti di vista diversi. Ogni \textit{stakeholder} guarda aspetti differenti:

\begin{figure}[H]
\centering
\begin{tikzpicture}
  \node[circle, fill=primary!12, draw=primary, line width=1.2pt, minimum size=2.6cm, font=\large\bfseries, text=primary!80!black] (q) at (0,0) {Qualità};
  \foreach \ang/\stake/\cosa/\col in {90/{Utente}/{«È facile da usare e veloce?»}/q4,
                                  18/{Cliente}/{«Fa ciò che ho chiesto, al prezzo concordato?»}/q1,
                                  -54/{Sviluppatore}/{«Il codice è leggibile e facile da modificare?»}/q7,
                                  -126/{Amministratore di sistema}/{«Si installa e convive con gli altri sistemi?»}/q8,
                                  162/{Società}/{«Protegge i dati? Può fare danni a persone?»}/q5}{
    \node[card=\col, text width=3.7cm, font=\scriptsize, align=center] (n\ang) at (\ang:4.1) {\textbf{\iuser\ \stake}\\\cosa};
    \draw[line width=1pt, \col] (q) -- (n\ang);
  }
\end{tikzpicture}
\caption{La qualità del software vista da stakeholder diversi.}
\end{figure}

Serve quindi un \textbf{modello condiviso} che metta ordine tra tutti questi aspetti. È ciò che forniscono gli standard internazionali.

% ============================================================
\section{I modelli di qualità ISO/IEC}

La famiglia di standard \textbf{ISO/IEC 25000} (detta \textit{SQuaRE}, \textit{Systems and software Quality Requirements and Evaluation})
si occupa della qualità del software. Lo standard \hi{ISO/IEC 25002} descrive i \textbf{modelli di qualità} del software, che sono due:

\begin{figure}[H]
\centering
\begin{tikzpicture}
  \node[card=swblue, text width=7.6cm, align=left, minimum height=4.3cm] (pq) {%
    \textbf{\large Product Quality Model}\\{\small\itshape Modello di qualità del prodotto}\\[4pt]\small
    Composto da \textbf{9 caratteristiche}, suddivise in sotto-caratteristiche, che riguardano le \textbf{proprietà del prodotto}.
    Fornisce un riferimento per \textbf{specificare, misurare e valutare} la qualità dei prodotti.\\[4pt]
    \scriptsize\textit{Punto di vista}: il software in sé.};
  \node[card=sworange, text width=7.6cm, align=left, minimum height=4.3cm, right=0.9cm of pq] (qu) {%
    \textbf{\large Quality-in-use Model}\\{\small\itshape Modello di qualità in uso}\\[4pt]\small
    Composto da \textbf{3 caratteristiche}, suddivise in sotto-caratteristiche, che descrivono l'effetto del prodotto sugli
    \textbf{stakeholder} quando viene \textbf{usato in uno specifico contesto d'uso}.\\[4pt]
    \scriptsize\textit{Punto di vista}: il software mentre qualcuno lo usa.};
  \draw[-{Stealth[length=3mm]}, line width=1.5pt, swgray] (pq.east) -- node[above, font=\scriptsize, align=center]{influenza} (qu.west);
\end{tikzpicture}
\caption{I due modelli di qualità: il primo guarda il prodotto, il secondo il prodotto in uso.}
\end{figure}

Entrambi i modelli sono \textbf{gerarchici}: la qualità si scompone in \textbf{caratteristiche}, ogni caratteristica in
\textbf{sotto-caratteristiche}, e queste ultime sono abbastanza concrete da poter essere \textbf{misurate}.

\begin{center}
\begin{tikzpicture}[font=\small]
  \node[chip=primary, minimum width=3cm, minimum height=0.7cm] (a) at (0,0) {Qualità};
  \node[chip=q2, minimum width=3cm, minimum height=0.7cm] (b) at (4.3,0) {Caratteristica};
  \node[chip=q2!60!white, text=black, minimum width=3.4cm, minimum height=0.7cm] (c) at (8.9,0) {Sotto-caratteristica};
  \node[chip=swgray, minimum width=3cm, minimum height=0.7cm] (d) at (13.4,0) {Misura};
  \draw[flow] (a) -- (b); \draw[flow] (b) -- (c); \draw[flow] (c) -- (d);
  \node[note, below=0.15cm of b, align=center] {es.\ Reliability};
  \node[note, below=0.15cm of c, align=center] {es.\ Availability};
  \node[note, below=0.15cm of d, align=center] {es.\ \% di tempo in cui\\il sistema è accessibile};
\end{tikzpicture}
\end{center}

\warning{Nota sui nomi}{
  Il modello del prodotto è descritto in dettaglio nello standard \textbf{ISO/IEC 25010}. Nella sua revisione più recente
  alcune caratteristiche hanno nomi diversi da quelli usati a lezione: ad esempio \textit{Usability} è chiamata
  \textit{Interaction Capability}. In questa dispensa manteniamo i nomi usati nelle slide del corso (in inglese, come nello standard).
}

% ============================================================
\section{Il modello di qualità del prodotto}

Il modello di qualità del prodotto è composto da nove caratteristiche:

\begin{figure}[H]
\centering
\begin{tikzpicture}
  \node[chip=primary, minimum width=6cm, minimum height=0.85cm, font=\small\bfseries] (root) at (0,0) {Software Product Quality};
  \foreach \i/\nome/\col in {0/{Functional\\Suitability}/q1, 1/{Reliability}/q2, 2/{Performance\\Efficiency}/q3,
                             3/{Usability}/q4, 4/{Security}/q5, 5/{Compat-\\ibility}/q6,
                             6/{Maintain-\\ability}/q7, 7/{Flexibility}/q8, 8/{Safety}/q9}{
    \node[chip=\col, minimum width=1.95cm, minimum height=1.0cm, font=\scriptsize\bfseries] (c\i) at ({(\i-4)*2.12},-1.9) {\nome};
    \draw[line width=0.9pt, swgray] (root.south) -- ++(0,-0.4) -| (c\i.north);
  }
\end{tikzpicture}
\caption{Le nove caratteristiche del modello di qualità del prodotto.}
\end{figure}

Vediamole una per una. Per ciascuna trovi una \textbf{domanda guida} che ne riassume il senso e una tabella delle sotto-caratteristiche
con un esempio concreto.

% ------------------------------------------------------------
\Needspace{12\baselineskip}\phantomsection\addcontentsline{toc}{subsection}{Functional Suitability}

\qualitybox{q1}{Functional Suitability}{Idoneità funzionale}{Fa quello che serve, e lo fa bene?}{
  Grado in cui un componente o un sistema fornisce funzioni che soddisfano le esigenze \textbf{dichiarate e implicite}
  quando viene usato nelle condizioni specificate.
}

\begin{table}[H]
\centering
\renewcommand{\arraystretch}{1.4}
\begin{tabularx}{\textwidth}{P{3.6cm} L L}
\rowcolor{q1!15}
\textbf{Sotto-caratteristica} & \textbf{Significato} & \textbf{Esempio (app di home banking)} \\
\textbf{Functional completeness}\newline\textit{\small completezza} & Le funzioni coprono \textbf{tutti} i compiti e gli obiettivi specificati. & L'app permette davvero di fare bonifici, vedere il saldo, pagare bollettini, come richiesto. \\
\textbf{Functional correctness}\newline\textit{\small correttezza} & Il prodotto fornisce risultati \textbf{accurati}. & Il saldo mostrato è esatto al centesimo. \\
\textbf{Functional appropriateness}\newline\textit{\small appropriatezza} & Le funzioni \textbf{facilitano} il raggiungimento dei compiti. & Per un bonifico a un contatto frequente bastano pochi passaggi, non dieci schermate. \\
\bottomrule
\end{tabularx}
\end{table}

% ------------------------------------------------------------
\Needspace{12\baselineskip}\phantomsection\addcontentsline{toc}{subsection}{Reliability}

\qualitybox{q2}{Reliability}{Affidabilità}{Funziona senza interruzioni, anche quando qualcosa va storto?}{
  Grado in cui il sistema svolge le sue funzioni nelle \textbf{condizioni specificate} per un \textbf{periodo di tempo specificato}.
}

\begin{table}[H]
\centering
\renewcommand{\arraystretch}{1.4}
\begin{tabularx}{\textwidth}{P{3.6cm} L L}
\rowcolor{q2!15}
\textbf{Sotto-caratteristica} & \textbf{Significato} & \textbf{Esempio} \\
\textbf{Faultlessness}\newline\textit{\small assenza di guasti} & Svolge le funzioni \textbf{senza guasti} durante il normale funzionamento. & L'app di messaggistica non si chiude inaspettatamente durante l'uso quotidiano. \\
\textbf{Availability}\newline\textit{\small disponibilità} & È \textbf{operativo e accessibile} quando serve. & Il sito delle prenotazioni è raggiungibile anche alle 8:00 del giorno di apertura degli appelli. \\
\textbf{Fault tolerance}\newline\textit{\small tolleranza ai guasti} & Funziona come previsto \textbf{nonostante} guasti hardware o software. & Se si rompe un disco del server, i dati restano accessibili grazie a una copia su un altro disco. \\
\textbf{Recoverability}\newline\textit{\small recuperabilità} & Dopo un'interruzione o un guasto, riesce a \textbf{ripristinarsi}. & Dopo un crash, l'editor riapre il documento con le modifiche salvate automaticamente. \\
\bottomrule
\end{tabularx}
\end{table}

La disponibilità si esprime spesso come \textbf{percentuale di tempo} in cui il sistema è accessibile. Le percentuali sembrano tutte
altissime, ma la differenza in termini di interruzioni è enorme:

\begin{table}[H]
\centering
\renewcommand{\arraystretch}{1.3}
\begin{tabular}{>{\bfseries}c c l}
\toprule
Disponibilità & Interruzione massima in un anno & In pratica\ldots \\
\midrule
99\%      & $\approx$ 3,65 giorni & Accettabile per un sito personale \\
99,9\%    & $\approx$ 8 ore e 46 minuti & Tipico di molti servizi web \\
99,99\%   & $\approx$ 52 minuti & Servizi importanti (e-commerce, banche) \\
99,999\%  & $\approx$ 5 minuti & Sistemi critici (telecomunicazioni, emergenze) \\
\bottomrule
\end{tabular}
\caption{Disponibilità e tempo di inattività annuo corrispondente (un anno = 8.760 ore).}
\end{table}

\info{{Guasto, errore, malfunzionamento}}{
  Quando si parla di affidabilità è utile distinguere tre concetti che nel linguaggio comune si confondono:
  \begin{center}
  \begin{tikzpicture}
    \node[card=q2, text width=3.9cm, font=\scriptsize, align=left, minimum height=1.9cm] (f) {\textbf{Difetto} (\textit{fault})\\Un errore nel codice o nel progetto.\\\textit{Es.}: un \texttt{<} scritto al posto di \texttt{<=}.};
    \node[card=sworange, text width=3.9cm, font=\scriptsize, align=left, minimum height=1.9cm, right=1.7cm of f] (e) {\textbf{Errore} (\textit{error})\\Uno stato interno sbagliato causato dal difetto.\\\textit{Es.}: il ciclo salta l'ultimo elemento.};
    \node[card=swred, text width=3.9cm, font=\scriptsize, align=left, minimum height=1.9cm, right=1.7cm of e] (m) {\textbf{Malfunzionamento} (\textit{failure})\\Comportamento sbagliato visibile all'esterno.\\\textit{Es.}: il totale della fattura è errato.};
    \draw[flow] (f) -- node[above, font=\scriptsize, align=center]{se\\eseguito} (e);
    \draw[flow] (e) -- node[above, font=\scriptsize, align=center]{se si\\propaga} (m);
  \end{tikzpicture}
  \end{center}
  Un difetto non sempre produce un malfunzionamento (magari quel codice non viene mai eseguito). La \textit{fault tolerance}
  consiste proprio nell'impedire che un difetto o un guasto diventi un malfunzionamento visibile.
}

% ------------------------------------------------------------
\Needspace{12\baselineskip}\phantomsection\addcontentsline{toc}{subsection}{Performance Efficiency}

\qualitybox{q3}{Performance Efficiency}{Efficienza prestazionale}{Risponde in fretta, senza sprecare risorse?}{
  Grado in cui un prodotto svolge le sue funzioni entro i \textbf{vincoli di risorse} specificati ed è efficiente nell'uso delle
  risorse (CPU, memoria, spazio di archiviazione, energia, \ldots).
}

\begin{table}[H]
\centering
\renewcommand{\arraystretch}{1.4}
\begin{tabularx}{\textwidth}{P{3.6cm} L L}
\rowcolor{q3!15}
\textbf{Sotto-caratteristica} & \textbf{Significato} & \textbf{Esempio} \\
\textbf{Time behaviour}\newline\textit{\small comportamento temporale} & Tempi di risposta e \textit{throughput} (lavoro svolto per unità di tempo) rispettano i requisiti. & La ricerca di un prodotto restituisce i risultati in meno di un secondo. \\
\textbf{Resource utilization}\newline\textit{\small uso delle risorse} & Quantità e tipo di risorse usate rispettano i requisiti. & L'app di navigazione non scarica la batteria dello smartphone in un'ora. \\
\textbf{Capacity}\newline\textit{\small capacità} & I \textbf{limiti massimi} di un parametro del sistema rispettano i requisiti. & Il sito di vendita dei biglietti regge 100.000 utenti connessi quando apre la vendita di un concerto. \\
\bottomrule
\end{tabularx}
\end{table}

% ------------------------------------------------------------
\Needspace{12\baselineskip}\phantomsection\addcontentsline{toc}{subsection}{Usability}

\qualitybox{q4}{Usability}{Usabilità}{Gli utenti riescono a usarlo senza fatica e senza sbagliare?}{
  Grado in cui un prodotto o sistema può essere \textbf{usato} dai suoi utenti, cioè quanto è facile interagire con esso.
}

\begin{table}[H]
\centering
\renewcommand{\arraystretch}{1.4}
\begin{tabularx}{\textwidth}{P{3.6cm} L L}
\rowcolor{q4!15}
\textbf{Sotto-caratteristica} & \textbf{Significato} & \textbf{Esempio} \\
\textbf{Recognizability}\newline\textit{\small riconoscibilità} & Gli utenti capiscono se il sistema è \textbf{adatto alle loro esigenze}. & La pagina iniziale spiega chiaramente a cosa serve il servizio. \\
\textbf{Learnability}\newline\textit{\small apprendibilità} & Le funzioni si imparano a usare \textbf{in un tempo specificato}. & Un nuovo cassiere impara a usare il gestionale del negozio in una mattina. \\
\textbf{Operability}\newline\textit{\small operabilità} & Il sistema è facile da \textbf{far funzionare e controllare}. & Le azioni più comuni sono raggiungibili con un clic. \\
\textbf{User error protection}\newline\textit{\small protezione dagli errori} & Il sistema \textbf{previene gli errori} dell'utente. & Prima di eliminare un file chiede conferma; per le date propone un calendario invece di un campo di testo libero. \\
\textbf{Inclusivity}\newline\textit{\small inclusività} & Può essere usato da persone con \textbf{background diversi} (età, abilità, culture, lingue, generi, situazioni economiche\ldots). & L'app è disponibile in più lingue, funziona con i lettori di schermo e resta usabile su smartphone economici. \\
\bottomrule
\end{tabularx}
\end{table}

L'usabilità avrà un ruolo centrale in una parte importante del corso (Modulo B). Può fare la differenza tra \textbf{successo e fallimento}
di un prodotto e, in alcuni contesti (dispositivi medici, cabine di pilotaggio, sale di controllo), perfino tra \textbf{la vita e la morte}.

\info{Prova tu: User Inyerface}{
  Il sito \url{https://userinyerface.com} è un «gioco» che simula un modulo di registrazione progettato apposta per essere
  \textbf{il più frustrante possibile}: pulsanti ingannevoli, caselle pre-selezionate al contrario, testi di aiuto che ostacolano.
  Provare a completarlo è il modo migliore per capire, per contrasto, cosa significa \textit{usabilità}.
}

% ------------------------------------------------------------
\Needspace{12\baselineskip}\phantomsection\addcontentsline{toc}{subsection}{Security}

\qualitybox{q5}{Security}{Sicurezza (protezione)}{I dati sono protetti da chi non dovrebbe accedervi o modificarli?}{
  Grado in cui un sistema si \textbf{difende dagli attacchi} di soggetti malintenzionati e protegge informazioni e dati, applicando
  adeguati meccanismi di \textbf{autorizzazione}.
}

\begin{table}[H]
\centering
\renewcommand{\arraystretch}{1.4}
\begin{tabularx}{\textwidth}{P{3.6cm} L L}
\rowcolor{q5!15}
\textbf{Sotto-caratteristica} & \textbf{Significato} & \textbf{Esempio (fascicolo sanitario / banca)} \\
\textbf{Confidentiality}\newline\textit{\small riservatezza} & I dati sono accessibili \textbf{solo a chi è autorizzato}. & Solo il paziente e i suoi medici possono leggere il referto. \\
\textbf{Integrity}\newline\textit{\small integrità} & Stato e dati sono protetti da \textbf{modifiche o cancellazioni non autorizzate}. & Nessuno, tranne il medico che l'ha redatto, può alterare il referto. \\
\textbf{Non-repudiation}\newline\textit{\small non ripudio} & Si può \textbf{dimostrare} che un'azione è avvenuta, così che non possa essere negata in seguito. & Un bonifico firmato digitalmente non può essere disconosciuto da chi l'ha disposto. \\
\textbf{Accountability}\newline\textit{\small responsabilità} & Le azioni di un'entità sono \textbf{tracciabili} in modo univoco fino a essa. & Un registro (\textit{log}) memorizza quale operatore ha consultato quale cartella clinica. \\
\textbf{Authenticity}\newline\textit{\small autenticità} & Si può \textbf{provare} che l'identità di un soggetto o di una risorsa è quella dichiarata. & L'accesso richiede credenziali più un codice inviato allo smartphone (autenticazione a due fattori). \\
\bottomrule
\end{tabularx}
\end{table}

% ------------------------------------------------------------
\Needspace{12\baselineskip}\phantomsection\addcontentsline{toc}{subsection}{Compatibility}

\qualitybox{q6}{Compatibility}{Compatibilità}{Convive e dialoga con gli altri sistemi?}{
  Grado in cui un sistema può \textbf{scambiare informazioni} con altri prodotti, sistemi o componenti, e/o svolgere le proprie funzioni
  \textbf{condividendo lo stesso ambiente} e le stesse risorse con altri sistemi.
}

\begin{table}[H]
\centering
\renewcommand{\arraystretch}{1.4}
\begin{tabularx}{\textwidth}{P{3.6cm} L L}
\rowcolor{q6!15}
\textbf{Sotto-caratteristica} & \textbf{Significato} & \textbf{Esempio} \\
\textbf{Co-existence}\newline\textit{\small coesistenza} & Funziona in modo efficiente condividendo ambiente e risorse con altri prodotti, \textbf{senza danneggiarli}. & L'antivirus non rallenta al punto da rendere inutilizzabili le altre applicazioni. \\
\textbf{Interoperability}\newline\textit{\small interoperabilità} & Può \textbf{scambiare informazioni} con altri sistemi e \textbf{usare} quelle ricevute. & Gli eventi dell'app universitaria si esportano in un formato standard e compaiono nel calendario dello smartphone. \\
\bottomrule
\end{tabularx}
\end{table}

% ------------------------------------------------------------
\Needspace{12\baselineskip}\phantomsection\addcontentsline{toc}{subsection}{Maintainability}

\qualitybox{q7}{Maintainability}{Manutenibilità}{Si può modificare facilmente, senza rompere nulla?}{
  Grado di efficacia ed efficienza con cui un prodotto può essere \textbf{modificato} per migliorarlo, correggerlo o adattarlo
  a cambiamenti dell'ambiente e dei requisiti.
}

\begin{table}[H]
\centering
\renewcommand{\arraystretch}{1.4}
\begin{tabularx}{\textwidth}{P{3.6cm} L L}
\rowcolor{q7!15}
\textbf{Sotto-caratteristica} & \textbf{Significato} & \textbf{Esempio} \\
\textbf{Modularity}\newline\textit{\small modularità} & Il software è composto da componenti distinti: una modifica a uno ha \textbf{impatto minimo} sugli altri. & Cambiare il fornitore dei pagamenti richiede di modificare solo il modulo \textit{Pagamenti}. \\
\textbf{Reusability}\newline\textit{\small riusabilità} & Un modulo può essere usato come risorsa in \textbf{più sistemi}. & Il modulo di autenticazione è usato sia dal sito sia dall'app mobile. \\
\textbf{Modifiability}\newline\textit{\small modificabilità} & Può essere modificato senza introdurre difetti o degradare la qualità. & Aggiungere l'IVA al calcolo del carrello richiede di toccare un solo metodo. \\
\textbf{Testability}\newline\textit{\small testabilità} & Si possono stabilire criteri di test ed \textbf{eseguire test} per verificarli. & Ogni classe può essere testata in isolamento con test automatici. \\
\bottomrule
\end{tabularx}
\end{table}

\info{Il collegamento con i costi}{
  Ricordando che la manutenzione rappresenta circa \textbf{due terzi} del costo di un software, la manutenibilità è una delle
  caratteristiche con il maggiore impatto economico. È esattamente ciò che si ottiene con una buona progettazione modulare
  (System Design) e con codice pulito (Implementation).
}

% ------------------------------------------------------------
\Needspace{12\baselineskip}\phantomsection\addcontentsline{toc}{subsection}{Flexibility}

\qualitybox{q8}{Flexibility}{Flessibilità}{Si adatta a nuovi ambienti, carichi e contesti?}{
  Grado in cui un prodotto può essere \textbf{adattato} a cambiamenti nei requisiti, nei contesti d'uso o nell'ambiente operativo.
}

\begin{table}[H]
\centering
\renewcommand{\arraystretch}{1.4}
\begin{tabularx}{\textwidth}{P{3.6cm} L L}
\rowcolor{q8!15}
\textbf{Sotto-caratteristica} & \textbf{Significato} & \textbf{Esempio} \\
\textbf{Adaptability}\newline\textit{\small adattabilità} & Può essere adattato o trasferito a \textbf{hardware, software o ambienti} diversi. & L'applicazione desktop funziona su Windows, macOS e Linux. \\
\textbf{Scalability}\newline\textit{\small scalabilità} & Gestisce carichi di lavoro \textbf{crescenti o decrescenti}, adattando la propria capacità. & Durante il Black Friday l'e-commerce aggiunge server automaticamente, e li rilascia a fine giornata. \\
\textbf{Installability}\newline\textit{\small installabilità} & Può essere \textbf{installato e disinstallato} con successo, in modo efficace ed efficiente. & L'installazione richiede pochi clic e la disinstallazione non lascia file sparsi nel sistema. \\
\textbf{Replaceability}\newline\textit{\small sostituibilità} & Può \textbf{sostituire} un altro prodotto con lo stesso scopo nello stesso ambiente. & Il nuovo gestionale importa i dati del vecchio, così l'azienda può passare da uno all'altro. \\
\bottomrule
\end{tabularx}
\end{table}

% ------------------------------------------------------------
\Needspace{12\baselineskip}\phantomsection\addcontentsline{toc}{subsection}{Safety}

\qualitybox{q9}{Safety}{Sicurezza (incolumità)}{Può fare del male a persone, beni o ambiente?}{
  Grado in cui un prodotto \textbf{evita} di trovarsi in uno stato in cui la vita umana, la salute, i beni o l'ambiente sono messi in pericolo.
}

\begin{table}[H]
\centering
\renewcommand{\arraystretch}{1.4}
\begin{tabularx}{\textwidth}{P{3.6cm} L L}
\rowcolor{q9!15}
\textbf{Sotto-caratteristica} & \textbf{Significato} & \textbf{Esempio} \\
\textbf{Fail safe}\newline\textit{\small sicurezza in caso di guasto} & In caso di guasto si porta \textbf{automaticamente} in una modalità operativa sicura. & Se il treno perde il segnale di controllo, frena automaticamente. \\
\textbf{Risk identification}\newline\textit{\small identificazione dei rischi} & Riconosce una sequenza di eventi o operazioni che può portare a un \textbf{rischio inaccettabile}. & Il microinfusore di insulina rileva che la dose richiesta è anomala rispetto allo storico. \\
\textbf{Hazard warning}\newline\textit{\small avviso di pericolo} & Avvisa dei rischi \textbf{in tempo utile} per poter reagire. & L'auto segnala il rischio di collisione con qualche secondo di anticipo. \\
\bottomrule
\end{tabularx}
\end{table}

\warning{Security e Safety: in italiano sono entrambe «sicurezza»}{
  In italiano usiamo una sola parola per due concetti ben distinti, e confonderli è un errore frequente:
  \begin{center}
  \begin{tikzpicture}
    % Security
    \node[chip=q5, minimum width=7.4cm, minimum height=0.7cm, font=\small\bfseries] (sh) at (0,0) {SECURITY: proteggere il sistema dal mondo};
    \node[card=swred, minimum width=1.9cm, font=\scriptsize] (att) at (-2.6,-1.35) {\iuser\\attaccante};
    \node[card=q5, minimum width=2.1cm, minimum height=1.1cm, font=\scriptsize\bfseries] (sys1) at (2.4,-1.35) {\icpu\\Sistema e dati};
    \draw[-{Stealth[length=3mm]}, line width=1.4pt, swred, decorate, decoration={zigzag, amplitude=1.5pt, segment length=6pt, post length=3mm}] (att) -- node[above=2pt, font=\tiny, text=black]{attacco} (sys1);
    % Safety
    \node[chip=q9, minimum width=7.4cm, minimum height=0.7cm, font=\small\bfseries] (sa) at (8.3,0) {SAFETY: proteggere il mondo dal sistema};
    \node[card=q9, minimum width=2.1cm, minimum height=1.1cm, font=\scriptsize\bfseries] (sys2) at (5.9,-1.35) {\icog\\Sistema};
    \node[card=swgreen, minimum width=2.4cm, font=\scriptsize] (ppl) at (10.8,-1.35) {\iuser\ \iuser\\persone, ambiente};
    \draw[-{Stealth[length=3mm]}, line width=1.4pt, q9, decorate, decoration={zigzag, amplitude=1.5pt, segment length=6pt, post length=3mm}] (sys2) -- node[above=2pt, font=\tiny, text=black]{danno} (ppl);
  \end{tikzpicture}
  \end{center}
  \textbf{Security} riguarda attacchi \textit{intenzionali} contro il sistema e i suoi dati (furto di password, accessi non autorizzati).
  \textbf{Safety} riguarda i \textit{danni} che il sistema può causare a persone, beni o ambiente, anche senza alcun attaccante
  (un difetto nel software di frenata di un'auto).
}

% ------------------------------------------------------------
\subsection{Il quadro completo}

\begin{tcbraster}[raster columns=3, raster equal height=rows, raster column skip=3mm, raster row skip=3mm]
  \begin{tcolorbox}[colback=q1!5, colframe=q1, title=\textbf{Functional Suitability}, fontupper=\small]
    \textbullet\ Functional Completeness\\\textbullet\ Functional Correctness\\\textbullet\ Functional Appropriateness
  \end{tcolorbox}
  \begin{tcolorbox}[colback=q2!5, colframe=q2, title=\textbf{Reliability}, fontupper=\small]
    \textbullet\ Faultlessness\\\textbullet\ Availability\\\textbullet\ Fault Tolerance\\\textbullet\ Recoverability
  \end{tcolorbox}
  \begin{tcolorbox}[colback=q3!5, colframe=q3, title=\textbf{Performance Efficiency}, fontupper=\small]
    \textbullet\ Time Behaviour\\\textbullet\ Resource Utilization\\\textbullet\ Capacity
  \end{tcolorbox}
  \begin{tcolorbox}[colback=q4!5, colframe=q4, title=\textbf{Usability}, fontupper=\small]
    \textbullet\ Recognizability\\\textbullet\ Learnability\\\textbullet\ Operability\\\textbullet\ User Error Protection\\\textbullet\ Inclusivity
  \end{tcolorbox}
  \begin{tcolorbox}[colback=q5!5, colframe=q5, title=\textbf{Security}, fontupper=\small]
    \textbullet\ Confidentiality\\\textbullet\ Integrity\\\textbullet\ Non-repudiation\\\textbullet\ Accountability\\\textbullet\ Authenticity
  \end{tcolorbox}
  \begin{tcolorbox}[colback=q6!5, colframe=q6, title=\textbf{Compatibility}, fontupper=\small]
    \textbullet\ Co-existence\\\textbullet\ Interoperability
  \end{tcolorbox}
  \begin{tcolorbox}[colback=q7!5, colframe=q7, title=\textbf{Maintainability}, fontupper=\small]
    \textbullet\ Modularity\\\textbullet\ Reusability\\\textbullet\ Modifiability\\\textbullet\ Testability
  \end{tcolorbox}
  \begin{tcolorbox}[colback=q8!5, colframe=q8, title=\textbf{Flexibility}, fontupper=\small]
    \textbullet\ Adaptability\\\textbullet\ Scalability\\\textbullet\ Installability\\\textbullet\ Replaceability
  \end{tcolorbox}
  \begin{tcolorbox}[colback=q9!5, colframe=q9, title=\textbf{Safety}, fontupper=\small]
    \textbullet\ Fail Safe\\\textbullet\ Risk Identification\\\textbullet\ Hazard Warning
  \end{tcolorbox}
\end{tcbraster}

% ============================================================
\section{Qualità, modularità e lavoro di squadra}

C'è un aspetto della qualità che si vede solo quando si passa dal «programma» al «prodotto software»: un software di qualità
deve essere \textbf{flessibile}, e per essere flessibile deve essere \textbf{modulare}. È il vecchio principio del
\hi{divide et impera}: dividere un problema grande in parti più piccole e gestibili.

La modularità non serve solo alla manutenzione: è ciò che rende possibile il \textbf{lavoro in gruppo}.

\begin{figure}[H]
\centering
\begin{tikzpicture}
  \node[card=swred, text width=6.4cm, minimum height=3.4cm, align=left] (mono) {%
    \textbf{Software monolitico}\\[4pt]\small
    Tutto è intrecciato: chi tocca una parte rischia di rompere le altre.\\[3pt]
    \textbullet\ difficile dividere il lavoro;\\
    \textbullet\ tutti devono conoscere tutto;\\
    \textbullet\ le modifiche si bloccano a vicenda.};
  \node[card=swgreen, text width=6.4cm, minimum height=3.4cm, align=left, right=3.2cm of mono] (modu) {%
    \textbf{Software modulare}\\[4pt]\small
    Il sistema è diviso in parti con responsabilità chiare.\\[3pt]
    \textbullet\ il lavoro si può \textbf{parallelizzare};\\
    \textbullet\ ogni gruppo lavora sul proprio modulo;\\
    \textbullet\ si sfruttano \textbf{competenze diverse}.};
  \draw[-{Stealth[length=3mm]}, line width=1.5pt, swgray] (mono) -- node[above, font=\scriptsize\bfseries, align=center]{divide\\et impera} (modu);
\end{tikzpicture}
\end{figure}

\begin{esempio}{Venti persone su un progetto}
Immaginiamo un progetto da realizzare con venti sviluppatori. Se il software è un unico blocco indistinto, venti persone
si ostacolano a vicenda: modificano gli stessi file, si sovrappongono, passano il tempo a coordinarsi.

Se invece il sistema è ben diviso in moduli, il lavoro si può distribuire e svolgere \textbf{in parallelo}: chi si occupa
dell'interfaccia grafica lavora sulle schermate, chi è esperto di basi di dati lavora sulla persistenza, chi conosce il dominio
implementa la logica applicativa. Ognuno lavora sul pezzo più adatto alle proprie competenze, e i moduli si integrano
attraverso interfacce concordate.
\end{esempio}

\info{Il legame con le caratteristiche ISO}{
  Questa proprietà non è un'opinione, ma è codificata nel modello di qualità: è la sotto-caratteristica \textbf{Modularity}
  (dentro \textit{Maintainability}), che a sua volta abilita \textbf{Reusability}, \textbf{Modifiability} e \textbf{Testability}.
  Ed è decisa molto presto, nella fase di \textbf{System Design}.
}

% ============================================================
\section{Il modello di qualità in uso}

Il secondo modello cambia prospettiva: non guarda il software «sul banco di prova», ma il software \textbf{mentre viene usato}.

\dfn{Qualità in uso}{
  Grado in cui un prodotto o sistema può essere usato da \textbf{specifici utenti} per soddisfare le loro esigenze e raggiungere
  \textbf{specifici obiettivi} con \textbf{efficacia}, \textbf{efficienza}, \textbf{assenza di rischi} e \textbf{soddisfazione},
  in \textbf{specifici contesti d'uso}.
}

\begin{figure}[H]
\centering
\begin{tikzpicture}
  \node[chip=primary, minimum width=5cm, minimum height=0.85cm, font=\small\bfseries] (root) at (0,0) {Quality-in-use};
  \foreach \k/\nome/\col/\subs/\x in {
      0/Usability/q4/{Effectiveness,Efficiency,Satisfaction}/-5.6,
      1/Safety/q9/{Commercial Damage,Operator Health and Safety,Public Health and Safety,Environmental Harm}/0,
      2/Flexibility/q8/{Context Conformity,Context Extensibility,Accessibility}/5.6}{
    \node[chip=\col, minimum width=4.2cm, minimum height=0.8cm, font=\small\bfseries] (c\k) at (\x,-1.6) {\nome};
    \draw[line width=0.9pt, swgray] (root.south) -- ++(0,-0.35) -| (c\k.north);
    \foreach \s [count=\n] in \subs {
      \node[subchip=\col, minimum width=3.6cm, minimum height=0.62cm, anchor=west] (s\k\n) at (\x-1.55,-1.6-\n*0.9) {\s};
      \draw[swgray, line width=0.7pt] (\x-1.85,-2.0) |- (s\k\n.west);
    }
  }
\end{tikzpicture}
\caption{Il modello di qualità in uso: 3 caratteristiche e le loro sotto-caratteristiche.}
\end{figure}

\subsection{Usability (in uso)}

Misura quanto gli utenti riescono a raggiungere i propri obiettivi usando il sistema in modo efficiente e soddisfacente.
\begin{itemize}
  \item \textbf{Effectiveness} (efficacia): quanto bene gli utenti \textbf{completano} i compiti che volevano svolgere.
  \item \textbf{Efficiency} (efficienza): quante \textbf{risorse} (tempo, fatica) servono per completarli.
  \item \textbf{Satisfaction} (soddisfazione): il \textbf{comfort} e l'esperienza positiva dell'utente durante l'uso.
\end{itemize}

\warning{Usability compare due volte: perché?}{
  \begin{itemize}
    \item Nel \textbf{modello del prodotto}, l'usabilità è misurata sulle \textbf{caratteristiche del prodotto}
          (l'interfaccia protegge dagli errori? è coerente? è accessibile?).
    \item Nel \textbf{modello in uso}, l'usabilità è misurata dal \textbf{punto di vista degli utenti} che lo usano davvero
          (sono riusciti a fare quello che volevano? in quanto tempo? erano soddisfatti?).
  \end{itemize}
  Esempio: un navigatore può avere pulsanti grandi e menu ben organizzati (usabilità del prodotto), ma se un automobilista
  di notte, sotto la pioggia, non riesce a impostare la destinazione, l'usabilità in uso in quel contesto è scarsa.
}

\subsection{Safety (in uso)}

Valuta la capacità del sistema di \textbf{prevenire danni} a persone, ambiente e interessi commerciali durante l'uso.
\begin{itemize}
  \item \textbf{Commercial damage}: quanto il sistema evita perdite economiche o danni all'attività (es.\ un errore nei prezzi dell'e-commerce).
  \item \textbf{Operator health and safety}: quanto protegge \textbf{chi lo usa} da rischi per la salute o la sicurezza.
  \item \textbf{Public health and safety}: quanto evita rischi o danni per il \textbf{pubblico} in generale.
  \item \textbf{Environmental harm}: quanto evita o riduce gli impatti negativi sull'\textbf{ambiente}.
\end{itemize}

\subsection{Flexibility (in uso)}

Indica la capacità del sistema di adattarsi e operare efficacemente in \textbf{contesti o ambienti diversi}.
\begin{itemize}
  \item \textbf{Context conformity}: capacità di adattarsi ai requisiti e ai vincoli specifici di contesti diversi.
  \item \textbf{Context extensibility}: potenzialità di estendersi o adattarsi a contesti nuovi o in cambiamento senza modifiche significative.
  \item \textbf{Accessibility}: quanto efficacemente il sistema può essere usato da \textbf{tutte le persone}, incluse quelle con disabilità, in ambienti diversi.
\end{itemize}

\subsection{Il legame tra i due modelli}

I due modelli non sono indipendenti: le proprietà del prodotto \textbf{influenzano} la qualità percepita durante l'uso,
ma quest'ultima dipende anche dal \textbf{contesto} (chi usa il sistema, per fare cosa, in quale ambiente).

\begin{figure}[H]
\centering
\begin{tikzpicture}
  \node[card=swblue, text width=4.8cm, minimum height=2.3cm, align=center] (p) {\textbf{Qualità del prodotto}\\[3pt]\small proprietà del software\\(misurabili sul prodotto)};
  \node[card=swgray, text width=3.6cm, minimum height=2.3cm, align=center, right=1.4cm of p] (ctx) {\textbf{Contesto d'uso}\\[3pt]\small utenti, obiettivi,\\ambiente fisico e sociale};
  \node[card=sworange, text width=4.8cm, minimum height=2.3cm, align=center, right=1.4cm of ctx] (u) {\textbf{Qualità in uso}\\[3pt]\small effetti su chi usa il sistema\\(efficacia, rischi, soddisfazione)};
  \draw[-{Stealth[length=3mm]}, line width=1.5pt, swblue] (p.north) to[out=30, in=150] node[above, font=\scriptsize]{influenza} (u.north);
  \draw[flow] (ctx) -- (u);
  \draw[-{Stealth[length=2.5mm]}, line width=1pt, dashed, swgray] (u.south) to[out=-150, in=-30] node[below, font=\scriptsize]{il feedback degli utenti guida i miglioramenti del prodotto} (p.south);
\end{tikzpicture}
\caption{Qualità del prodotto e qualità in uso: la seconda dipende dalla prima e dal contesto.}
\end{figure}

% ============================================================
\section{Qualità e compromessi}

Un'ultima osservazione fondamentale: le caratteristiche di qualità \textbf{non si possono massimizzare tutte insieme}.
Spesso migliorarne una ne peggiora un'altra.

\begin{table}[H]
\centering
\renewcommand{\arraystretch}{1.45}
\begin{tabularx}{\textwidth}{P{2.9cm} >{\centering\arraybackslash}p{0.6cm} P{2.9cm} L}
\toprule
\multicolumn{3}{c}{\textbf{Caratteristiche in conflitto}} & \textbf{Perché} \\
\midrule
\textcolor{q5}{\textbf{Security}} & $\leftrightarrow$ & \textcolor{q4}{\textbf{Usability}} & Password complesse e autenticazione a due fattori proteggono meglio, ma rendono l'accesso più scomodo. \\
\textcolor{q3}{\textbf{Performance}} & $\leftrightarrow$ & \textcolor{q7}{\textbf{Maintainability}} & Il codice di Quake III è velocissimo ma quasi impossibile da capire e modificare. \\
\textcolor{q9}{\textbf{Safety}} & $\leftrightarrow$ & \textcolor{q2}{\textbf{Reliability}} (availability) & Un treno che frena a ogni anomalia è sicuro, ma il servizio si interrompe spesso. \\
\textcolor{q8}{\textbf{Flexibility}} & $\leftrightarrow$ & \textcolor{q3}{\textbf{Performance}} & Un'applicazione che gira su ogni piattaforma rinuncia spesso alle ottimizzazioni specifiche di ciascuna. \\
\bottomrule
\end{tabularx}
\caption{Alcuni compromessi tipici tra caratteristiche di qualità.}
\end{table}

\info{Il ruolo dell'ingegnere}{
  Per questo i requisiti di qualità vanno \textbf{esplicitati e ordinati per priorità} già in fase di Requirements Engineering:
  per il software di un pacemaker la \textit{safety} viene prima di tutto, per un videogioco contano soprattutto prestazioni e soddisfazione.
  Trovare il giusto equilibrio, dati il problema e i vincoli, è esattamente il lavoro dell'ingegnere del software.
}

% ============================================================
\section{Esercizi}

\begin{esercizio}{Quale caratteristica di qualità?}
Per ciascun requisito indica la caratteristica e la sotto-caratteristica del modello di qualità del prodotto a cui si riferisce.
\begin{enumerate}
  \item «Il portale deve essere raggiungibile almeno il 99,9\% del tempo.»
  \item «La pagina dei risultati di ricerca deve essere visualizzata entro 2 secondi.»
  \item «Solo il medico curante può leggere la cartella clinica del paziente.»
  \item «In caso di perdita del segnale, il drone deve atterrare automaticamente.»
  \item «Un nuovo utente deve riuscire a effettuare un ordine senza consultare il manuale entro 10 minuti dal primo accesso.»
  \item «L'app deve importare i contatti esportati da altre applicazioni in formato standard.»
  \item «Il sistema deve reggere un traffico dieci volte superiore durante i saldi, riducendo le risorse quando il picco finisce.»
  \item «Ogni operazione sul conto deve essere registrata insieme all'identità dell'operatore che l'ha eseguita.»
  \item «Il modulo di spedizione deve poter essere sostituito senza modificare gli altri moduli.»
\end{enumerate}

\Needspace{14\baselineskip}
\textbf{Soluzione.}
\begin{center}
\renewcommand{\arraystretch}{1.3}
\begin{tabular}{c l l}
\toprule
\textbf{n.} & \textbf{Caratteristica} & \textbf{Sotto-caratteristica} \\
\midrule
1 & Reliability & Availability \\
2 & Performance Efficiency & Time behaviour \\
3 & Security & Confidentiality \\
4 & Safety & Fail safe \\
5 & Usability & Learnability \\
6 & Compatibility & Interoperability \\
7 & Flexibility & Scalability \\
8 & Security & Accountability \\
9 & Maintainability & Modularity \\
\bottomrule
\end{tabular}
\end{center}
\end{esercizio}

\begin{esercizio}{Security o Safety?}
\begin{enumerate}
  \item Un malintenzionato prova migliaia di password per entrare nell'account di un utente.
  \item Il software di un robot industriale lo ferma se un operatore entra nella sua area di lavoro.
  \item I voti degli studenti non devono poter essere modificati da chi non è docente.
  \item Il sistema di controllo di una diga deve avvisare in tempo in caso di livello dell'acqua pericoloso.
\end{enumerate}

\textbf{Soluzione.}
(1) Security (authenticity: il sistema deve verificare l'identità).
(2) Safety (fail safe / operator health and safety).
(3) Security (integrity).
(4) Safety (hazard warning).
\end{esercizio}

\gbox{In sintesi}{azure-gradient-6!80!black}{
\begin{itemize}
  \item La qualità del software \textbf{non è una sola}: dipende dal prodotto e dal punto di vista degli stakeholder, ed è difficile da quantificare.
  \item Gli standard \textbf{ISO/IEC 25000} (SQuaRE) definiscono due modelli gerarchici: \textbf{qualità del prodotto} e \textbf{qualità in uso}.
  \item Il modello del prodotto ha \textbf{9 caratteristiche}: Functional Suitability, Reliability, Performance Efficiency, Usability,
        Security, Compatibility, Maintainability, Flexibility, Safety.
  \item Il modello in uso ha \textbf{3 caratteristiche}: Usability (efficacia, efficienza, soddisfazione), Safety e Flexibility.
  \item \textbf{Security} protegge il sistema dagli attacchi; \textbf{Safety} protegge persone e ambiente dal sistema.
  \item La \textbf{modularità} (divide et impera) è ciò che rende il software flessibile e permette di dividere il lavoro in un team.
  \item Le caratteristiche sono spesso in \textbf{conflitto}: vanno prioritizzate in base al contesto.
\end{itemize}
}
